% Options for packages loaded elsewhere
\PassOptionsToPackage{unicode}{hyperref}
\PassOptionsToPackage{hyphens}{url}
%
\documentclass[
]{article}
\usepackage{amsmath,amssymb}
\usepackage{lmodern}
\usepackage{iftex}
\ifPDFTeX
  \usepackage[T1]{fontenc}
  \usepackage[utf8]{inputenc}
  \usepackage{textcomp} % provide euro and other symbols
\else % if luatex or xetex
  \usepackage{unicode-math}
  \defaultfontfeatures{Scale=MatchLowercase}
  \defaultfontfeatures[\rmfamily]{Ligatures=TeX,Scale=1}
\fi
% Use upquote if available, for straight quotes in verbatim environments
\IfFileExists{upquote.sty}{\usepackage{upquote}}{}
\IfFileExists{microtype.sty}{% use microtype if available
  \usepackage[]{microtype}
  \UseMicrotypeSet[protrusion]{basicmath} % disable protrusion for tt fonts
}{}
\makeatletter
\@ifundefined{KOMAClassName}{% if non-KOMA class
  \IfFileExists{parskip.sty}{%
    \usepackage{parskip}
  }{% else
    \setlength{\parindent}{0pt}
    \setlength{\parskip}{6pt plus 2pt minus 1pt}}
}{% if KOMA class
  \KOMAoptions{parskip=half}}
\makeatother
\usepackage{xcolor}
\IfFileExists{xurl.sty}{\usepackage{xurl}}{} % add URL line breaks if available
\IfFileExists{bookmark.sty}{\usepackage{bookmark}}{\usepackage{hyperref}}
\hypersetup{
  pdftitle={Gastroenteritis-Erreger am Beispiel E.Coli},
  pdfauthor={Jakob Schumacher},
  hidelinks,
  pdfcreator={LaTeX via pandoc}}
\urlstyle{same} % disable monospaced font for URLs
\usepackage[margin=1in]{geometry}
\usepackage{graphicx}
\makeatletter
\def\maxwidth{\ifdim\Gin@nat@width>\linewidth\linewidth\else\Gin@nat@width\fi}
\def\maxheight{\ifdim\Gin@nat@height>\textheight\textheight\else\Gin@nat@height\fi}
\makeatother
% Scale images if necessary, so that they will not overflow the page
% margins by default, and it is still possible to overwrite the defaults
% using explicit options in \includegraphics[width, height, ...]{}
\setkeys{Gin}{width=\maxwidth,height=\maxheight,keepaspectratio}
% Set default figure placement to htbp
\makeatletter
\def\fps@figure{htbp}
\makeatother
\setlength{\emergencystretch}{3em} % prevent overfull lines
\providecommand{\tightlist}{%
  \setlength{\itemsep}{0pt}\setlength{\parskip}{0pt}}
\setcounter{secnumdepth}{-\maxdimen} % remove section numbering
\ifLuaTeX
  \usepackage{selnolig}  % disable illegal ligatures
\fi

\title{Gastroenteritis-Erreger am Beispiel E.Coli}
\author{Jakob Schumacher}
\date{08. April 2022}

\begin{document}
\maketitle

\hypertarget{gastroenteritis-erreger-am-beispiel-e.coli}{%
\section{Gastroenteritis-Erreger am Beispiel
E.Coli}\label{gastroenteritis-erreger-am-beispiel-e.coli}}

Jakob Schumacher

\hypertarget{uxfcber-den-vortrag}{%
\subsection{Über den Vortrag}\label{uxfcber-den-vortrag}}

Dies ist ein Vortrag von Jakob Schumacher für den Kurs: WBKD-01 Akademie
für öffentliches Gesundheitswesen . Die Zielgruppe sind Amtsärztinnen
und Amtsärzte. Der Vortrag wurde am 08.04.2022 gehalten.

Online

PDF

Word

\includegraphics{../assets/2022-04-08-WBKD-Gastroenteritis/QR-link.svg}

\includegraphics{../assets/2022-04-08-WBKD-Gastroenteritis/QR-pdf.svg}

\includegraphics{../assets/2022-04-08-WBKD-Gastroenteritis/QR-doc.svg}

Der Vortrag wurde erstellt mit \href{https://revealjs.com/}{Reveal.js}
in \href{https://www.rstudio.com/}{R/RStudio}. Die Präsentationswebseite
läuft mit \href{https://jekyllrb.com/}{Jekyll}.

Download als PDF

\hypertarget{uxfcber-mich}{%
\subsection{Über mich}\label{uxfcber-mich}}

\begin{itemize}
\tightlist
\item
  Hygienereferent im Gesundheitsamt Reinickendorf
\item
  Zuvor:

  \begin{itemize}
  \tightlist
  \item
    Robert Koch-Institut
  \item
    Bundesministerium für Gesundheit
  \item
    Medizin (Krankenhaus, Gefängnis)
  \end{itemize}
\item
  Projekte:

  \begin{itemize}
  \tightlist
  \item
    \href{https://jakobschumacher.github.io/infektionsschutzkompendium/}{Infektionsschutzkompendium}
  \item
    \href{http://merkblaetter-fuer-infektionskrankheiten.de/}{Merkblätter
    für Infektionskrankheiten}
  \end{itemize}
\end{itemize}

\hypertarget{teil-1-rolle-der-gesundheitsuxe4mter}{%
\section{Teil 1: Rolle der
Gesundheitsämter}\label{teil-1-rolle-der-gesundheitsuxe4mter}}

\hypertarget{malerische-beschreibung}{%
\subsection{Malerische Beschreibung}\label{malerische-beschreibung}}

\hypertarget{aufgabe-der-gesundheitsuxe4mter-im-infektionsschutz}{%
\subsection{Aufgabe der Gesundheitsämter im
Infektionsschutz}\label{aufgabe-der-gesundheitsuxe4mter-im-infektionsschutz}}

\hypertarget{aufgaben-der-amtsleitunginfektionsschutzleitung}{%
\subsection{Aufgaben der
Amtsleitung/Infektionsschutzleitung}\label{aufgaben-der-amtsleitunginfektionsschutzleitung}}

\begin{itemize}
\tightlist
\item
  Methoden erarbeiten/trainieren
\item
  Maßnahmen hochhalten
\item
  Ausbruchserkennung organisieren
\item
  Krisenvorbereitung
\item
  Netzwerk pflegen
\item
  Hintergrundwissen Mikrobiologie
\item
  Personal aquirieren/schulen
\end{itemize}

\hypertarget{teil-2-uxfcberblick}{%
\section{Teil 2 Überblick}\label{teil-2-uxfcberblick}}

\hypertarget{uxfcberblilck-uxfcber-gastroenteritis-erreger}{%
\subsection{Überblilck über Gastroenteritis
Erreger}\label{uxfcberblilck-uxfcber-gastroenteritis-erreger}}

\begin{itemize}
\tightlist
\item
  Norovirus
\item
  Rotavirus
\item
  Salmonellen
\item
  Campylobacter
\item
  Shigellen
\item
  E.Coli
\item
  Yersiniose
\item
  Listeriose
\item
  Giardia
\item
  Kryptosporidiose
\end{itemize}

\hypertarget{die-wichtigsten-punkte}{%
\subsection{Die wichtigsten Punkte}\label{die-wichtigsten-punkte}}

\begin{itemize}
\tightlist
\item
  Gestroenteritis-Erreger sind häufig
\item
  Durchfall ist ``gemeinsame Endstrecke'' von Viren, Bakterien,
  Parasiten
\item
  Großer Killer der Welt: 1,3 Mio (bei 2 Milliarden Infizierten)
\item
  Gastroenteritis-Erreger Vorbeugung: Wassermanagement, Toiletten,
  Landwirtschaft
\item
  In Deutschland größte Gefahr durch: EHEC/HUS
\item
  Bei der Ermittlung: Ausbrüche bedenken
\item
  Zusammenarbeit mit Lebensmittelaufsicht
\item
  Norovirus ist ein Indikator für Hygiene
\end{itemize}

\hypertarget{der-darm}{%
\subsection{Der Darm}\label{der-darm}}

\begin{itemize}
\tightlist
\item
  4-5 Meter lang
\item
  Enthält 1-2 Kilogramm Bakterien (Mikrobiom)
\item
  Wichtigste Aufgabe: Verdauung, Immunabwehr und ?
\item
  Das Mikrobiom spielt eine größere Rolle als bislang gedacht
\end{itemize}

\hypertarget{was-muss-man-zu-einem-erreger-wissen}{%
\subsection{Was muss man zu einem Erreger
wissen}\label{was-muss-man-zu-einem-erreger-wissen}}

\begin{enumerate}
\def\labelenumi{\arabic{enumi}.}
\tightlist
\item
  Namen des Erregers und Krankheit
\item
  Epidemiologie
\item
  Vorkommen
\item
  Umweltverhalten
\item
  Übertragungswege
\item
  Zeiten
\item
  Labornachweise
\item
  Meldepflichten
\item
  Infektionsinterview-Fragen
\item
  Maßnahmen
\end{enumerate}

\hypertarget{teil-3}{%
\section{Teil 3}\label{teil-3}}

E.Coli und seine Krankheiten

\hypertarget{uxfcberblick-uxfcber-e.coli}{%
\subsection{Überblick über E.Coli}\label{uxfcberblick-uxfcber-e.coli}}

\begin{itemize}
\tightlist
\item
  Coliforme Erreger

  \begin{itemize}
  \tightlist
  \item
    E.Coli

    \begin{itemize}
    \tightlist
    \item
      ETEC
    \item
      EPEC
    \item
      EIEC
    \item
      EAEC
    \item
      EHEC
    \item
      MRGN
    \end{itemize}
  \end{itemize}
\item
  Je nach Toxin unterschiedliche Krankheitsbilder
\end{itemize}

\hypertarget{e.coli-als-symbiont}{%
\subsection{E.Coli als Symbiont}\label{e.coli-als-symbiont}}

\begin{itemize}
\tightlist
\item
  Hilft mit bei der Nahrungsaufspaltung
\item
  Produziert Vitamine (B-Vitamine, K-Vitamin)
\end{itemize}

\hypertarget{e.coli-als-pathogen}{%
\subsection{E.Coli als Pathogen}\label{e.coli-als-pathogen}}

\begin{itemize}
\tightlist
\item
  Gastroenteritis
\item
  Infektionen
\item
  HUS
\item
  (Indikatorkeim)
\end{itemize}

\hypertarget{gastroenteritis-symptome}{%
\subsection{Gastroenteritis-Symptome}\label{gastroenteritis-symptome}}

\begin{itemize}
\tightlist
\item
  Krankheitsgefühl
\item
  Bauchschmerzen
\item
  Durchfall
\item
  Erbrechen
\item
  Schmerzen
\item
  Erkältungssymptome (Bei Erkältung auch Gastroenteritis)
\item
  Fieber (Zeichen für Immunaktivität)
\item
  Blut im Stuhl (Zeichen für Invastion)
\item
  Exsikkose
\end{itemize}

\hypertarget{huxe4molytisch-uruxe4mische-symptome-hus}{%
\subsection{Hämolytisch-Urämische Symptome
(HUS)}\label{huxe4molytisch-uruxe4mische-symptome-hus}}

\begin{itemize}
\tightlist
\item
  Anämie (Blutarmut)
\item
  Urämie (Nierenversagen)
\item
  Thrombozytopenie (Blutung)
\item
  Neurologische Symptome
\end{itemize}

\hypertarget{infektionen}{%
\subsection{Infektionen}\label{infektionen}}

\begin{itemize}
\tightlist
\item
  Wundinfektion
\item
  Harnwegsinfektionen
\item
  Sepsis
\end{itemize}

\hypertarget{e.coli-und-die-toxine}{%
\subsection{E.Coli und die Toxine}\label{e.coli-und-die-toxine}}

\begin{itemize}
\tightlist
\item
  E.Coli hat diverse Mechanismen/Toxine
\item
  Manche Mechanismen schützen vor Antibiotika
\item
  Manche Mechanismen verursachen eine besser Anheftung oder Eindringung
  in die Zelle
\item
  Shigatoxin 2 (nur EHEC) verursacht ein HUS
\end{itemize}

\hypertarget{epidemiologie}{%
\subsection{Epidemiologie}\label{epidemiologie}}

\begin{itemize}
\tightlist
\item
  Massive Inzidenzen in Ländern mit hygienisch niedrigen Standards
\item
  Deutschland

  \begin{itemize}
  \tightlist
  \item
    \textasciitilde10.000 E.Coli-Enteritis
  \item
    \textasciitilde1000 Ehec
  \item
    \textasciitilde100 HUS
  \end{itemize}
\item
  Kinder und Reiserückkehrer
\end{itemize}

\hypertarget{vorkommen}{%
\subsection{Vorkommen}\label{vorkommen}}

\begin{itemize}
\tightlist
\item
  E.Coli lebt im Enddarm von vielen Warmblütern
\item
  Weltweit vertreten
\item
  EHEC vor allem in Wiederkäuern: Rinder, Schafe, Ziegen
\end{itemize}

\hypertarget{umweltverhalten}{%
\subsection{Umweltverhalten}\label{umweltverhalten}}

\begin{itemize}
\tightlist
\item
  E.Coli lebt im Darm von Tieren
\item
  Von dort auf das Fell und die Umgebung
\item
  Vom Fell beim Schlachten aufs Fleisch
\item
  Vom Fell auf die Hand
\item
  Vom Darm auf die Felder
\item
  Vom Darm ins Wasser
\item
  Vom Wasser aufs Feld
\end{itemize}

\hypertarget{wichtigste-uxfcbertragung-auf-menschen}{%
\subsection{Wichtigste Übertragung auf
Menschen}\label{wichtigste-uxfcbertragung-auf-menschen}}

\begin{itemize}
\tightlist
\item
  Kontaktinfektion (Schmierinfektion, Fäkal-Oral, Händeinfektion)
\item
  Lebensmittel
\item
  Trinkwasser
\end{itemize}

\hypertarget{zeiten}{%
\subsection{Zeiten}\label{zeiten}}

\begin{itemize}
\tightlist
\item
  Inkubationszeit: 2-10 Tage
\item
  Ansteckungsfähige Zeit: Solange Symptome bestehen. Aber
  Dauerausscheidung möglich
\end{itemize}

\hypertarget{labornachweise}{%
\subsection{Labornachweise}\label{labornachweise}}

\begin{itemize}
\tightlist
\item
  Kultur (Goldstandard)
\item
  PCR (Weißt die Gene nach)
\item
  ELISA (Weißt die Toxine nach)
\item
  Genomsequenzierung
\end{itemize}

\hypertarget{e.coli-erscheint-als-folgende-meldung}{%
\subsection{E.Coli erscheint als folgende
Meldung}\label{e.coli-erscheint-als-folgende-meldung}}

\begin{itemize}
\tightlist
\item
  HUS (Arzt)
\item
  EHEC (Labor)
\item
  E.Coli-Enteritis (Labor,
  \href{https://www.gesetze-im-internet.de/ifsg/__34.html}{§34},
  \href{https://www.gesetze-im-internet.de/ifsg/__36.html}{§36})
\item
  Wasserkeim (TrinkWV)
\item
  MRGN (Labor)
\item
  Wundinfektionskeim (Fehlmeldung)
\end{itemize}

\hypertarget{mauxdfnahmen-ehec-i}{%
\subsection{Maßnahmen EHEC I}\label{mauxdfnahmen-ehec-i}}

\begin{itemize}
\tightlist
\item
  Informationsbeschaffung, Übermittlung, Informationsweitergabe
  (Lebensmittelaufsicht)
\item
  Ausbruchsuntersuchung
\item
  Aufklärung

  \begin{itemize}
  \tightlist
  \item
    Erreger
  \item
    Küchenhygiene
  \item
    Toilettenhygiene.
  \item
    Händehygiene
  \item
    Wäschehygiene
  \item
    Lebensmittelhygiene auf Reisen
  \item
    Streichelzoo
  \end{itemize}
\end{itemize}

\hypertarget{mauxdfnahmen-ehec-ii}{%
\subsection{Maßnahmen EHEC II}\label{mauxdfnahmen-ehec-ii}}

\begin{itemize}
\tightlist
\item
  Isolierung der erkrankten Person in einer Einrichtung
\item
  Tätigkeitsverbot nach § 34
\item
  Tätigkeitsverbot nach § 42
\end{itemize}

\hypertarget{infektionsinterview-fragen-i}{%
\subsection{Infektionsinterview-Fragen
I}\label{infektionsinterview-fragen-i}}

\begin{itemize}
\tightlist
\item
  Wie geht es Ihnen? (!)
\item
  Welche Krankheitszeichen sind bei Ihnen aufgetreten?
\item
  Wann hat die Erkrankung genau begonnen?
\item
  Sind in Ihrem Umfeld weitere Personen erkrankt?
\item
  Haben Sie eine Verbindung zu einer Kita oder ähnlichen Einrichtung?
\item
  Haben Sie eine Verbindung zu einer Einrichtung, in der sehr viele
  Menschen untergebracht sind?
\end{itemize}

\hypertarget{infektionsinterview-fragen-ii}{%
\subsection{Infektionsinterview-Fragen
II}\label{infektionsinterview-fragen-ii}}

\begin{itemize}
\tightlist
\item
  Wohnen Sie mit jemand zusammen, der in eine Kita oder Schule geht oder
  dort arbeitet?
\item
  Haben Sie beruflich mit Lebensmitteln zu tun?
\item
  Waren Sie vor der Erkrankung im Ausland?
\item
  Waren sie vor Erkrankungsbeginn im Kontakt mit Wiederkäuern?
\end{itemize}

\hypertarget{teil-4-gruppenarbeit}{%
\section{Teil 4: Gruppenarbeit}\label{teil-4-gruppenarbeit}}

\hypertarget{teil-5}{%
\section{Teil 5}\label{teil-5}}

\hypertarget{wichtige-links}{%
\section{Wichtige Links}\label{wichtige-links}}

\begin{itemize}
\tightlist
\item
  Infektionsschutzkompendium
\item
  RKI-Ratgeber
\item
  Lebensmittelhygiene BFR
\item
  Reisewarnungen des Auswärtigen Amtes
\end{itemize}

\end{document}
